% POKE-FISI — Informe del proyecto (plantilla ACL, versión final en español)
% Compilar con pdfLaTeX:  pdflatex -> bibtex -> pdflatex -> pdflatex
% Si babel-spanish diera problemas de compilación, comenta la línea de \usepackage{babel}.
\documentclass[11pt]{article}

\usepackage{acl}            % versión final (sin la opción [review])
\usepackage{times}
\usepackage{latexsym}
\usepackage[T1]{fontenc}
\usepackage[utf8]{inputenc}
\usepackage{microtype}
\usepackage{inconsolata}
\usepackage{graphicx}
\usepackage[spanish,es-tabla]{babel}   % tipografía y "Tabla" (no "Cuadro")

\setlength\titlebox{6.5cm}   % más espacio para el título largo + 3 autores

\title{Agentes de Inteligencia Artificial para Combate Pokémon por Turnos: de Heurísticas a Minimax con Pesos Optimizados por Algoritmos Genéticos}

% TODO: reemplazar "Segundo/Tercer/Cuarto Autor" y sus correos por los reales.
% Misma institución para los cuatro: nombres en fila, afiliación compartida y
% los correos en dos filas (cuerpo pequeño) para que no se desborden.
\author{
  \textbf{Omar Cavero} \quad \textbf{Segundo Autor} \quad
  \textbf{Tercer Autor} \quad \textbf{Cuarto Autor} \\
  Universidad Nacional Mayor de San Marcos \\
  Facultad de Ingeniería de Sistemas e Informática \\[2pt]
  {\small
  \texttt{omar.cavero@unmsm.edu.pe} \quad
  \texttt{segundo.autor@unmsm.edu.pe} \\
  \texttt{tercer.autor@unmsm.edu.pe} \quad
  \texttt{cuarto.autor@unmsm.edu.pe}}
}

\begin{document}
\maketitle

\begin{abstract}
Presentamos POKE-FISI, un simulador de combate Pokémon por turnos diseñado para
estudiar agentes de inteligencia artificial de complejidad creciente. Sobre un
motor de combate determinista en lo estructural y estocástico en el daño,
implementamos tres niveles de agente: (1) un agente aleatorio de referencia,
(2) una heurística voraz de un turno basada en la diferencia de puntos de vida, y
(3) un agente Minimax con poda $\alpha$-$\beta$ cuya función de evaluación combina
linealmente cinco componentes, con pesos optimizados mediante un algoritmo
genético. Tratamos la naturaleza simultánea del juego con un colapso secuencial
``paranoico'' y empleamos números aleatorios comunes para comparaciones de baja
varianza y reproducibles. Nuestros experimentos muestran que el agente Minimax
supera al voraz, pero que el valor proviene casi en su totalidad de la
\emph{búsqueda} (que se satura a profundidad dos) y no del ajuste de pesos: tanto
la convergencia del algoritmo genético como un análisis de ablación coinciden en
que solo dos componentes (puntos de vida y número de unidades vivas) aportan
señal, mientras que ventaja de tipo, velocidad y presión de derribo son
marginales frente a un rival voraz. Reportamos este resultado negativo de forma
honesta y discutimos sus causas.
\end{abstract}

\section{Introducción}

Los juegos por turnos de información (casi) perfecta son un banco de pruebas
clásico para la búsqueda adversarial y la optimización de funciones de
evaluación \citep{shannon1950,russell2010}. En este trabajo construimos
POKE-FISI, un simulador de combates tipo Pokémon, y lo usamos para comparar
agentes de inteligencia artificial de dificultad creciente, desde una línea base
aleatoria hasta un agente Minimax cuyos pesos se optimizan con un algoritmo
genético (AG).

El problema tiene características que lo hacen interesante: los equipos y los
movimientos se asignan aleatoriamente al inicio de cada combate, ambos
jugadores eligen su acción de forma \emph{simultánea}, y existe la opción de
cambiar de unidad durante la partida. Esto introduce decisiones estratégicas
(cuándo atacar y cuándo cambiar) que una heurística miope no captura.

Nuestras contribuciones son:
\begin{itemize}
  \item Un motor de combate puro (sin dependencias gráficas) que permite simular
        miles de partidas de forma reproducible.
  \item Tres agentes de complejidad creciente bajo una interfaz común,
        incluyendo un Minimax con poda $\alpha$-$\beta$ que maneja el movimiento
        simultáneo y los cambios de unidad.
  \item Una función de evaluación de cinco componentes cuyos pesos se optimizan
        con un AG, junto con una metodología experimental reproducible basada en
        números aleatorios comunes.
  \item Un análisis honesto que, mediante la convergencia del AG y un estudio de
        ablación, identifica qué componentes realmente importan y muestra que,
        en este dominio, el ajuste de pesos aporta un margen marginal sobre la
        búsqueda.
\end{itemize}

\section{Trabajo relacionado}

La búsqueda adversarial Minimax y la poda $\alpha$-$\beta$ son herramientas
fundacionales del juego por computadora \citep{shannon1950,knuth1975}, llevadas a
su máxima expresión en sistemas como Deep Blue \citep{campbell2002}. La calidad de
un agente Minimax depende en gran medida de su función de evaluación de las hojas,
cuya construcción manual es laboriosa. Una alternativa es \emph{aprender} los
pesos de una evaluación lineal mediante computación evolutiva
\citep{holland1975,goldberg1989,eiben2015}. Nuestro trabajo combina ambas ideas en
un dominio de movimiento simultáneo y, a diferencia de gran parte de la
literatura centrada en resultados positivos, documenta con rigor el alcance
(limitado) de la optimización de pesos frente a un rival débil.

\section{Modelado del combate}
\label{sec:modelo}

\paragraph{Reglas.} Cada entrenador dispone de un equipo de 3 o 4 unidades, con
una activa por bando. En cada turno ambos jugadores eligen simultáneamente una
acción: atacar con uno de sus cuatro movimientos o cambiar la unidad activa. Los
cambios se resuelven antes que los ataques; entre dos ataques, actúa primero el de
mayor prioridad y, en empate, el de mayor velocidad. Cuando una unidad llega a
cero puntos de vida (PV) se debilita y su entrenador debe enviar un reemplazo. El
combate termina cuando un bando se queda sin unidades.

\paragraph{Fórmula de daño.} El daño base sigue la expresión del curso,
\begin{equation}
  \label{eq:dano}
  D = \left(\frac{A}{D_{\text{op}}}\right)\cdot BP \; - \; V_{\text{op}}\cdot K,
\end{equation}
donde $A$ es el ataque (físico o especial) del atacante, $D_{\text{op}}$ la
defensa correspondiente del oponente, $BP$ la potencia del movimiento,
$V_{\text{op}}$ la velocidad del oponente y $K{=}0.35$ una constante. Sobre
$D$ se aplican los bonificadores estándar: bonificación por tipo coincidente
(STAB, $\times 1.5$), efectividad de tipo $E$ (tabla de 18 tipos, generación 6+),
golpe crítico ($\times 1.5$ con probabilidad $1/16$), un factor aleatorio
uniforme en $[0.85, 1.0]$ y una tirada de precisión.

\paragraph{Datos.} El juego incluye 36 unidades con estadísticas reales
escaladas a nivel 50 y 144 movimientos (18 tipos $\times$ 8) con tipo, categoría,
potencia, precisión y puntos de poder. Los datos viven en archivos JSON
editables, de modo que el rebalanceo no requiere tocar código.

\section{Agentes}
\label{sec:agentes}

Todos los agentes implementan una interfaz común que, dado un estado y un bando,
devuelve una acción. Esto permite añadir niveles nuevos sin reescribir lo
existente.

\subsection{Nivel 1: Aleatorio}
El agente de referencia elige uniformemente entre las acciones legales del turno.
No distingue entre atacar y cambiar. Sirve de línea base: cualquier agente
posterior debe superarlo holgadamente.

\subsection{Nivel 2: Heurística voraz}
El agente de Nivel 2 realiza una búsqueda voraz de un solo turno (\emph{lookahead}
1-ply). Para cada acción legal simula un turno asumiendo que el rival usa su
movimiento de mayor daño esperado, y puntúa el estado resultante con la heurística
\begin{equation}
  \label{eq:h1}
  H_1 = \frac{\sum PV_{\text{prop}}}{\sum PV^{\max}_{\text{prop}}}
      - \frac{\sum PV_{\text{riv}}}{\sum PV^{\max}_{\text{riv}}} \in [-1,1],
\end{equation}
es decir, la diferencia normalizada de PV de equipo. Dentro del \emph{lookahead}
el daño se calcula sin azar (promediando el factor aleatorio y multiplicando por
la precisión), lo que hace al agente determinista y a los experimentos
reproducibles.

\subsection{Nivel 3: Minimax con poda $\alpha$-$\beta$}
\label{sec:n3}

\paragraph{Movimiento simultáneo.} El Minimax clásico supone turnos alternados,
pero en este juego ambos jugadores deciden a la vez. Adoptamos un \emph{colapso
secuencial paranoico}: suponemos que elegimos primero y que el rival responde
óptimamente \emph{viendo} nuestra jugada. Al concederle más información de la que
posee, obtenemos una estimación pesimista y segura, e integrable directamente con
poda $\alpha$-$\beta$.

\paragraph{Transición determinista.} La búsqueda no usa el motor estocástico, sino
una transición propia con \emph{daño esperado}
\begin{equation}
  \label{eq:dhat}
  \hat{D} = D \cdot \text{STAB} \cdot E \cdot \bar{r} \cdot \text{acc},
\end{equation}
con $\bar{r}{=}0.925$ (promedio del factor aleatorio), sin golpe crítico y con la
precisión como multiplicador. Así el árbol es determinista. Un punto clave de
implementación es que la clonación de estados para la búsqueda emplea un generador
aleatorio independiente, de modo que la deliberación del agente \emph{no} altera
el azar real del combate; esto es imprescindible para que los números aleatorios
comunes (Sección~\ref{sec:metodo}) funcionen durante el entrenamiento.

\paragraph{Función de evaluación.} Las hojas se puntúan con una combinación
lineal
\begin{equation}
  \label{eq:eval}
  H(s) = \sum_{i=1}^{5} w_i\, C_i(s),
\end{equation}
donde cada componente $C_i \in [-1,1]$ es un diferencial ``propio menos rival'':
(C1) diferencia normalizada de PV; (C2) diferencia normalizada de unidades vivas;
(C3) ventaja de tipo de la unidad activa, en escala logarítmica simétrica;
(C4) ventaja de velocidad de la activa; y (C5) presión de derribo (si puedo
debilitar al rival este turno menos si él puede debilitarme a mí). Los estados
terminales reciben un valor que domina cualquier valor heurístico, con una
bonificación por ganar antes.

\paragraph{Búsqueda.} Implementamos Minimax con poda $\alpha$-$\beta$. Para mejorar
la poda, ordenamos los ataques por daño esperado descendente y limitamos los
cambios a los $K$ mejores por calidad de emparejamiento (poda dirigida de
cambios). Los cambios forzados (tras un debilitamiento) se resuelven en solitario,
replicando la dinámica real del juego.

\paragraph{Optimización de pesos con un AG.} Los pesos $w_i$ se optimizan con un
algoritmo genético. El genoma es el vector $\mathbf{w}$ restringido al símplex
($w_i \ge 0$, $\sum_i w_i = 1$), lo que fija la escala y aporta interpretabilidad.
La aptitud (\emph{fitness}) de un genoma es su tasa de victorias contra el
Nivel~2, más términos pequeños de margen de PV y rapidez de victoria para densar
la señal. Usamos selección por torneo con elitismo, cruce BLX-$\alpha$, mutación
gaussiana con $\sigma$ decreciente y reparación al símplex. La selección final se
hace sobre un conjunto \emph{held-out} de configuraciones nunca vistas en
entrenamiento.

\section{Metodología experimental}
\label{sec:metodo}

Todos los experimentos enfrentan al agente evaluado contra el Nivel~2 sobre
equipos y movimientos generados aleatoriamente. Para reducir varianza usamos
\textbf{números aleatorios comunes} (CRN): todas las condiciones (o todos los
genomas de una generación) juegan exactamente las mismas configuraciones, lo que
convierte la comparación en pareada. Las configuraciones se re-muestrean entre
generaciones para evitar sobreajuste, y la siembra de los generadores aleatorios
no depende del orden de ejecución, de modo que el entrenamiento es reproducible y
paralelizable sin alterar los resultados. Reportamos tasa de victorias (con su
error estándar), margen de PV final y duración media en turnos. Salvo indicación,
los combates son 4 contra 4 y la profundidad de búsqueda es 2.

\section{Resultados}
\label{sec:resultados}

Como referencia, el Nivel~2 vence al Nivel~1 en torno al 82\% de las partidas,
confirmando que la heurística voraz mejora claramente sobre el azar.

\paragraph{Efecto de la profundidad.} La Tabla~\ref{tab:depth} muestra que la
búsqueda aporta valor al pasar de profundidad 1 a 2, pero \emph{se satura}: la
profundidad 3 no mejora y multiplica el costo. El desajuste entre el modelo
paranoico (rival óptimo) y el rival real (voraz) ayuda a explicar esta saturación.

\begin{table}[t]
  \centering
  \begin{tabular}{lc}
    \hline
    \textbf{Profundidad} & \textbf{Victorias N3} \\
    \hline
    1 & 55\% \\
    2 & 62\% \\
    3 & 62\% \\
    \hline
  \end{tabular}
  \caption{Barrido de profundidad del Minimax contra el Nivel~2 (100 partidas,
  CRN, pesos del prior). La profundidad satura en 2.}
  \label{tab:depth}
\end{table}

\paragraph{Comparación de pesos.} La Tabla~\ref{tab:pesos} compara distintos
vectores de pesos sobre un conjunto independiente de 300 partidas. Los pesos
convergidos del AG obtienen la mejor tasa (59.0\%) y el mejor margen, pero todos
los vectores ``sensatos'' quedan en $\sim$57--59\%, indistinguibles entre sí
dentro del error ($\pm 2.9\%$). El único claramente inferior es el uniforme
(54.0\%), que malgasta peso en componentes poco útiles.

\begin{table}[t]
  \centering
  \begin{tabular}{lcc}
    \hline
    \textbf{Pesos} & \textbf{Victorias} & \textbf{Margen PV} \\
    \hline
    uniforme                 & 54.0\% & $+0.076$ \\
    prior (manual)           & 57.3\% & $+0.137$ \\
    solo PV                  & 57.7\% & $+0.141$ \\
    macro (PV+vivos)         & 58.7\% & $+0.149$ \\
    \textbf{AG convergido}   & \textbf{59.0\%} & \textbf{$+0.156$} \\
    \hline
  \end{tabular}
  \caption{Comparación de vectores de pesos contra el Nivel~2
  (300 partidas, CRN, $\pm 2.9\%$). Diferencias menores a $\sim 6\%$ no son
  concluyentes.}
  \label{tab:pesos}
\end{table}

\paragraph{Convergencia del AG.} Entrenado con 25 individuos, 60 partidas por
evaluación y 50 generaciones, el AG convirgió de forma estable a pesos dominados
por PV y unidades vivas ($\approx 90\%$ del peso combinado), llevando velocidad y
presión de derribo a $\approx 0$.

\paragraph{Ablación.} La Tabla~\ref{tab:ablacion} cuantifica la contribución de
cada componente anulándolo desde un vector uniforme. Solo PV y unidades vivas
aportan señal estadísticamente significativa ($\sim 6\%$ cada uno); tipo,
velocidad y presión de derribo son indistinguibles de cero. Esto coincide
exactamente con el punto al que convergió el AG.

\begin{table}[t]
  \centering
  \begin{tabular}{lc}
    \hline
    \textbf{Componente anulado} & \textbf{Contribución} \\
    \hline
    Unidades vivas          & $+6.0\%$ \\
    Puntos de vida          & $+5.7\%$ \\
    Ventaja de tipo         & $+1.3\%$ \\
    Velocidad               & $+0.3\%$ \\
    Presión de derribo (KO) & $+0.3\%$ \\
    \hline
  \end{tabular}
  \caption{Ablación \emph{leave-one-out} desde pesos uniformes
  (300 partidas, $\pm 2.9\%$). Solo PV y unidades vivas son significativos.}
  \label{tab:ablacion}
\end{table}

\section{Discusión}
\label{sec:discusion}

Los resultados convergen en una conclusión clara y, en parte, contraintuitiva.
Primero, \textbf{la búsqueda aporta el grueso del valor}: el Minimax supera al
voraz, pero la ventaja se agota a profundidad 2. Segundo, \textbf{el ajuste de
pesos tiene un margen marginal} contra el Nivel~2: todos los vectores razonables
rinden $\sim$58\%, y el AG no supera a un prior hecho a mano. Tercero, y más
informativo, \textbf{dos vías independientes} (la convergencia del AG y la
ablación) coinciden en que solo PV y unidades vivas importan en este dominio.

¿Por qué? El juego es mecánicamente simple, por lo que un agente voraz de un turno
ya está cerca de lo óptimo en muchas posiciones; el azar de equipos domina
numerosas partidas con independencia de la habilidad; y el modelo paranoico
sobreestima a un rival que en realidad juega de forma voraz. En conjunto, el techo
de victorias contra el Nivel~2 ronda el 58\%.

Un aprendizaje metodológico relevante es la \emph{maldición del ganador}: una
selección \emph{held-out} con pocas partidas (80) eligió un candidato que parecía
alcanzar 62.5\%, pero que en un conjunto fresco mayor regresó a $\sim$56\%. Por
ello, optamos por shippear los pesos \emph{convergidos} del AG (representativos y
mejor medidos) en lugar del ``mejor de \emph{held-out}''.

El valor real del AG en este proyecto, por tanto, no es batir un win-rate sino
\emph{descubrir automáticamente} la estructura de la evaluación (qué componentes
importan), coincidiendo con una ablación manual. Es un resultado negativo, pero
riguroso y reproducible.

\section{Conclusiones}
\label{sec:conclusiones}

Presentamos POKE-FISI y una jerarquía de agentes que culmina en un Minimax con
poda $\alpha$-$\beta$ y pesos optimizados por un AG. El agente supera a la
heurística voraz, pero mostramos honestamente que, en este dominio y frente a un
rival voraz, la ganancia procede de la búsqueda (saturada en profundidad 2) y no
del ajuste fino de pesos, de los cuales solo dos componentes (PV y unidades vivas)
aportan señal. Como trabajo futuro, evaluar el AG contra un rival más fuerte
(por ejemplo, en juego propio o contra un Minimax de mayor profundidad) podría dar
a la optimización de pesos un margen genuino donde demostrar su valor.

\section*{Limitaciones}

El estudio se circunscribe a un motor de combate simplificado: los movimientos de
estado no tienen efecto, no hay condiciones alteradas, clima ni habilidades, y la
fórmula de daño es la del curso, no la de la serie. El rival de referencia
(Nivel~2) es débil, lo que satura la métrica de win-rate y limita el margen
observable de la optimización de pesos; las conclusiones sobre la (ir)relevancia
de ciertos componentes son específicas de este rival y podrían cambiar frente a
oponentes más fuertes. Finalmente, los tamaños muestrales (100--300 partidas)
implican errores estándar de varios puntos porcentuales, por lo que diferencias
pequeñas no son concluyentes.

\section*{Agradecimientos}

Agradecemos al curso de Inteligencia Artificial de la Universidad Nacional Mayor
de San Marcos. Los recursos artísticos y de audio empleados en la interfaz
provienen de fuentes de uso académico no comercial.

\bibliography{referencias}

\appendix

\section{Detalles de reproducibilidad}
\label{sec:appendix}

El entrenamiento del AG usa una semilla maestra fija de la que se derivan los
generadores de la batalla y de cada agente como flujos independientes. Los
hiperparámetros reportados son: población 25, 60 partidas por evaluación, 50
generaciones, profundidad de entrenamiento 2, poda de cambios top-2, torneo de
tamaño 3 con 2 élites, cruce BLX-0.3 y mutación gaussiana con $\sigma_0{=}0.25$ y
decaimiento 0.8. La evaluación final se realiza a profundidad 2--3 sobre
conjuntos \emph{held-out}. El entrenamiento es paralelizable a nivel de genoma y
produce resultados idénticos a la ejecución en serie.

\end{document}
